\newcommand{\INFMIN}{0.914}
\newcommand{\INFMAX}{0.995}
\newcommand{\CROSSMAX}{0.460}
\newcommand{\CROSSMIN}{0.049}
\newcommand{\CROSSMAXAVG}{0.232}
\newcommand{\SPMLPMIN}{0.11}
\newcommand{\SPMLPMAX}{0.24}
\newcommand{\KSMED}{0.48}
\newcommand{\KSSHARE}{45}

\newcommand{\GITHUB}{\href{https://github.com/officialteekay/cs412-nids-generalization}{\nolinkurl{github.com/officialteekay/cs412-nids-generalization}}}

\newcommand{\ABSTRACT}{Machine learning intrusion detectors often report F1 scores above 0.95, but those scores usually come from training and testing on the same dataset. This project tests what happens when a detector trained on one network is moved to a different one. Logistic regression, a random forest, a multilayer perceptron, and a small self-attention model were trained on each of three NetFlow v2 datasets, UNSW-NB15, ToN-IoT, and CSE-CIC-IDS2018, and tested on all three. Every model scores between \INFMIN{} and \INFMAX{} F1 on its own network. On networks it was not trained on, its average F1 falls to between \CROSSMIN{} and \CROSSMAXAVG{}, which is below the 0.357 a detector gets by flagging every flow. The random forest is the best model in-dataset and the worst across networks, with an average cross-dataset AUC of 0.413. Training on two networks does not help with the third. The models rely on different features on each network, and in UNSW-NB15 a single rule on the TTL field reaches an F1 of 0.832 while flagging every flow on the other networks. Removing the TTL features does not fix transfer, because the models switch to other network-specific shortcuts.}

\newcommand{\ATTNCROSS}{The attention model, trained on the 60,000 flow subsample, reaches 0.232.}

\newcommand{\INTRORESULTS}{The results are clear. Every model looks strong in-dataset, with F1 between \INFMIN{} and \INFMAX{}. On networks they were not trained on, the average F1 of every model falls to between \CROSSMIN{} and \CROSSMAXAVG{}, which is lower than the 0.357 a detector gets by labeling every flow as an attack. The random forest is the best model in-dataset and the worst across networks, and its average cross-dataset AUC of 0.413 is below random guessing. The attention model gets a higher cross-dataset F1 than an MLP trained on the same data, but its AUC is the same and the extra F1 comes from flagging more flows. Training on two networks does not help with a third. The feature analysis shows that the models depend on different features on each network, and that a one-line rule on the TTL field reaches an F1 of 0.832 on UNSW-NB15. Removing the TTL features does not make the models transfer better.}

\newcommand{\ATTNPARA}{The self-attention model is the most complex model in this project, so it is the strongest test of this idea. It was trained on 60,000 flows per source, so the fair comparison is MLP-60k, which saw exactly the same training flows. Shrinking the training set barely changed the MLP. Its in-dataset F1 went from 0.968 to 0.964 and its cross-dataset F1 from 0.157 to 0.146, so the amount of training data is not what limits transfer.

Attention-60k has a higher average cross-dataset F1 than MLP-60k, 0.232 compared with 0.146, and its gap of 0.730 is the smallest of the five models. That looks like a real improvement at first, but the other numbers do not support it. Its average cross-dataset AUC is 0.496, almost the same as the 0.495 of MLP-60k, so it is no better at separating attacks from benign flows on a new network. Its cross-dataset false positive rate is higher, 0.398 compared with 0.322. Most of the F1 gain comes from the two cells where ToN-IoT is the test set. There, 64 percent of the flows are attacks, so a model that flags more flows gets a higher F1 almost automatically. The attention model also varies the most between seeds, with a standard deviation of 0.049 in its gap compared with 0.002 for MLP-60k. Overall, the attention model does not show a real advantage across networks.}

\newcommand{\FLAGPARA}{Table~\ref{tab:summary} also includes a trivial detector that labels every flow as an attack. Its F1 is not zero, because F1 depends on how common attacks are in the test set. It scores 0.077 on UNSW-NB15, 0.214 on CIC-IDS2018, and 0.780 on ToN-IoT, for an average of 0.357. Every trained model has a lower average cross-dataset F1 than this trivial detector. Cell by cell, logistic regression beats it in 2 of the 6 cross-dataset pairs, the MLP, MLP-60k, and attention model in 1 each, and the random forest in none. Flagging every flow is useless in practice because its false positive rate is 100 percent, but the comparison shows how little the cross-dataset F1 scores mean on their own.}

\newcommand{\MULTICAPTION}{Logistic regression, the random forest, and the MLP use three seeds. MLP-60k and Attention-60k use one seed and 30,000 flows from each source.}

\newcommand{\MULTIATTN}{MLP-60k, trained on the same flows as the attention model, goes from 0.146 to 0.134. The attention model is the only exception, going from 0.232 to 0.306. Most of that gain is on ToN-IoT, where it reaches 0.608, and its held-out AUC also rises to 0.663 on ToN-IoT and 0.806 on CIC-IDS2018. On CIC-IDS2018, though, it gets its F1 of 0.297 by flagging 57.3 percent of the benign flows, and on UNSW-NB15 its AUC drops to 0.182. This result comes from a single seed, and the attention model was the least stable model across seeds in the single-source runs, so it is weak evidence at best.}

\newcommand{\ABLATIONPARA}{If TTL were the main reason the UNSW-NB15 models fail elsewhere, removing it should help. Table~\ref{tab:ablation} shows that it does not. Without MIN\_TTL and MAX\_TTL, the in-dataset F1 on UNSW-NB15 barely changes for the random forest, 0.952 to 0.951, and the MLP, 0.940 to 0.940. Only logistic regression drops, from 0.914 to 0.869. This means UNSW-NB15 has other features that separate the classes just as well, so TTL was not the only shortcut available.

Cross-dataset F1 got worse instead of better. For the models trained on UNSW-NB15, average F1 on the other two networks fell from 0.227 to 0.174 for logistic regression, 0.088 to 0.023 for the random forest, and 0.263 to 0.091 for the MLP. The reason is in the recall. With TTL, a UNSW-NB15 model flags a large share of the other networks just because their TTL values are not 31. The MLP trained on UNSW-NB15 flagged 69.0 percent of the attacks in CIC-IDS2018 this way, but also 46.8 percent of the benign flows. Without TTL, that accidental flagging stops. Its recall on CIC-IDS2018 falls to 2.4 percent and its false positive rate to 15.7 percent. Average cross-dataset AUC did not improve for any model, 0.614 to 0.580 for logistic regression, 0.413 to 0.413 for the random forest, and 0.493 to 0.473 for the MLP. Removing one artifact moved the failure around but did not make the models learn anything that transfers.}

\newcommand{\CONCLABLATION}{Removing the TTL features did not help, because the models found other shortcuts on each network.}
